\chapter{RVV and the pystencils RVV Backend}\label{r02:ch:rvv-primer}

This chapter gives the minimum of RISC-V Vector background that the rest of the report
needs, explains how \ps{} vectorizes a kernel, and shows what the RVV backend makes of
the waLBerla LBM sweep. A longer primer is \file{docs/notes/rvv-isa-primer.md}; the
backend itself is documented in Report~1, \cref{r01:ch:rvv-backend}.

\section{RVV in five concepts}\label{r02:sec:rvv-concepts}

The ratified RISC-V \emph{V} extension, version 1.0~\cite{rvv-spec}, adds 32 vector
registers \code{v0}--\code{v31} and instructions on them. Unlike SSE, AVX or NEON, the
ISA does not fix the register width.
\begin{description}
  \item[VLEN] The width of one vector register in bits, chosen by the hardware vendor
    (a power of two, at least 128 for the full V extension). The SpacemiT~X60 cores
    of the K1 have \textbf{VLEN = 256}. A program can read it at run time
    (\code{csrr vlenb}). The extension \code{Zvl256b} in a \code{-march} string promises
    the compiler that VLEN $\geq$ 256; the compiler then defines
    \code{__riscv_v_min_vlen} as 256.
  \item[SEW] The selected element width: 8, 16, 32 or 64 bits.
  \item[LMUL] The register-group multiplier. With LMUL = 2, 4 or 8 an instruction works
    on 2, 4 or 8 registers at once; fractional values (\code{mf2}, \code{mf4},
    \code{mf8}) use part of one register.
  \item[VLMAX] The number of elements an instruction can process:
    $\text{VLMAX} = \text{LMUL}\cdot\text{VLEN}/\text{SEW}$. For \code{double} at LMUL 1
    on the K1: $1\cdot256/64 = 4$.
  \item[vl] The \term{active vector length}, $\leq$ VLMAX, set by \code{vsetvli}. Every
    vector instruction processes only the first \code{vl} elements.
\end{description}
In C and C++, RVV is programmed through intrinsics~\cite{rvv-intrinsics}. Their names
encode the operation, the operand kinds and the type, and every intrinsic takes the
active length as its last argument:
\begin{ccode}
vfloat64m1_t a = __riscv_vle64_v_f64m1(&x[i], 4);      // load 4 doubles (f64, LMUL 1)
vfloat64m1_t c = __riscv_vfadd_vv_f64m1(a, b, 4);      // vector + vector
__riscv_vse64_v_f64m1(&y[i], c, 4);                    // store
\end{ccode}
The vector types (\code{vfloat64m1_t}, \code{vint32mf2_t}, \dots) are \term{sizeless}:
their size is unknown at compile time, so they cannot be members of a struct or
elements of an array. RVV provides \term{tuple types} such as \code{vfloat32m1x4_t}
instead.

\section{Vector-length agnostic or specific}\label{r02:sec:vla-vls}

Because VLEN varies between chips, RVV code is usually written \term{vector-length
agnostic} (VLA): a stripmining loop asks the hardware in every iteration how many
elements it may process, and the same binary runs on any VLEN. The alternative is
\term{vector-length specific} (VLS) code: assume one VLEN, and treat RVV like a
fixed-width SIMD instruction set.
\begin{center}\small
\begin{tabularx}{\linewidth}{@{}l>{\raggedright\arraybackslash}X>{\raggedright\arraybackslash}X@{}}
  \toprule
  & \textbf{VLA (stripmining)} & \textbf{VLS (fixed width)}\\
  \midrule
  loop & \code{for (; n > 0; n -= vl) { vl = vsetvl(n); ... }} & vector loop over
    $\lfloor n/L\rfloor L$ elements, scalar remainder\\
  portability & any VLEN & only VLEN $\geq$ the assumed value (and exactly equal for
    full efficiency)\\
  remainder & handled by the last, shorter \code{vl} & separate scalar loop\\
  code generator & needs new loop structures and IR & same structure as AVX/NEON\\
  \bottomrule
\end{tabularx}
\end{center}
The backend used here is \textbf{VLS with VLEN = 256}. \Cref{r02:sec:design-vls}
explains why.

\section{How pystencils vectorizes}\label{r02:sec:ps-vectorize}

\ps{} vectorizes the innermost loop of a kernel. With
\code{cpu.vectorize.enable = True}, the code generator:
\begin{enumerate}
  \item splits the innermost ($x$) loop into a vector loop that advances by $L$ lanes
        and a scalar remainder loop (\cref{r02:fig:fzyx});
  \item rewrites the loop body into vector IR: vector types, vector memory accesses,
        broadcasts of loop-invariant scalars, horizontal reductions;
  \item hands the vector IR to the \term{platform} of the target, whose
        \term{intrinsic selector} (a subclass of \cls{SelectIntrinsics}) replaces every
        vector operation by the target's intrinsic: \code{_mm256_add_pd} on AVX,
        \code{vaddq_f64} on NEON, \code{__riscv_vfadd_vv_f64m1} on RVV.
\end{enumerate}
The number of lanes $L$ comes from \code{Target.default_vector_lanes(dtype)}, e.g.\ 4
for \code{double} on AVX. Contiguous vector loads need the innermost stride to be 1,
which the user promises with \code{cpu.vectorize.assume_inner_stride_one}. Without
that promise, accesses along $x$ must use strided loads (on RVV) or gathers.

\section{The RVV backend at a glance}\label{r02:sec:backend-glance}

The backend added in \ps{} (\cref{r02:ch:pystencils}) follows the contract of the AVX
backend line by line:
\begin{center}\small
\begin{tabularx}{\linewidth}{@{}>{\raggedright\arraybackslash}p{4.5cm}>{\raggedright\arraybackslash}X@{}}
  \toprule
  \textbf{AVX (x86)} & \textbf{RVV backend (VLEN = 256)}\\
  \midrule
  \code{__m256d}, 4 doubles; \code{__m256}, 8 floats & \code{vfloat64m1_t},
    \code{vfloat32m1_t}: one register (LMUL 1)\\
  \code{__m128} (half width) & LMUL \code{mf2}, only for elements $\leq$ 32 bits\\
  no wider vectors & no register groups (no \code{m2}, \code{m4})\\
  implicit length & literal lane count as \code{vl} in every call (\code{4uLL}), which
    equals VLMAX at VLEN = 256\\
  \code{_mm256_set_pd}, horizontal helpers in \file{x86.hpp} & \code{set_<sfx>},
    \code{horizontal_<op>_<sfx>} in \file{rvv.hpp}\\
  \code{-march=x86-64-v3} & \code{-march=rv64gcv_zvl256b}\\
  \bottomrule
\end{tabularx}
\end{center}

\Cref{r02:lst:kernel} shows the generated kernel of the waLBerla stream-collide sweep
(\code{explicit_cse} variant, double precision) as it was produced on the K1. Apart from
reformatting, elisions and the \code{//} comments, it is unchanged.

\rneedspace{12\baselineskip}
\lstinputlisting[style=cpp,nolol=false,basicstyle=\ttfamily\scriptsize,
  caption={[Generated RVV kernel of the D3Q19 stream-collide sweep]Generated RVV kernel of the D3Q19 stream-collide sweep (reformatted and abridged, comments added in \texttt{//}).},
  label=r02:lst:kernel]{streamcollide-kernel.cpp}

Things to notice:
\begin{itemize}
  \item The structure is exactly that of an AVX kernel: hoisted constants, the $z$ and
        $y$ loops, a vector loop over $x$ in steps of 4, and a scalar remainder.
  \item Every access to the PDF field is a unit-stride \code{vle64}/\code{vse64}, which
        is what \code{assume_inner_stride_one} buys. The neighbour offsets of the pull
        stream are just constant offsets in the index expressions.
  \item Per vector iteration the source contains 70 loads (from 19 distinct
        addresses), 23 stores (19 PDFs, $\rho$ and 3 velocity components), 101
        multiplications, 82 additions, 43 subtractions and one division ($1/\rho$).
        Constants are broadcast with \code{vfmv_v_f}.
  \item RVV intrinsics are opaque to the compilers' algebraic optimizations. Neither GCC
        nor Clang contracts an intrinsic multiply and add into a fused multiply-add,
        or folds a multiplication by $-1$; on x86 the intrinsics are ordinary vector
        operations in the headers, so this never mattered there. GCC also does not
        merge repeated identical intrinsic calls. This is why the global common
        subexpression elimination of lbmpy (\code{cse_global}) matters for GCC
        (\cref{r02:sec:results-single}).
\end{itemize}
